\chapter{Theoretical Background}
\label{sec:background}

% Key concepts, definitions, and theoretical frameworks.
% Review of relevant literature; establish the state of research.

\section{Competitive advantage and firm performance}
\label{sec:constructs}

Competitive advantage is measured in 15 of the 67 studies reviewed here, and the
remaining 52 report performance outcomes only. Barney separates the two constructs by
what rivals are able to do: a firm holds a competitive advantage while it runs a
value-creating strategy no rival, actual or potential, is running at the same time, and
that advantage counts as sustained only once those rivals also cannot reproduce what the
strategy yields \parencite[p.~102]{barney1991firm}\footnote{[Barney (1991), p.~102] ``In
this article, a firm is said to have a \emph{competitive advantage} when it is
implementing a value creating strategy not simultaneously being implemented by any current
or potential competitors. A firm is said to have a \emph{sustained competitive advantage}
when it is implementing a value creating strategy not simultaneously being implemented by
any current or potential competitors \emph{and} when these other firms are unable to
duplicate the benefits of this strategy.''}. Sustained in that sense says nothing about
elapsed time. It asks whether the advantage is still there after rivals have given up
trying to copy it \parencite[p.~102]{barney1991firm}\footnote{[Barney (1991), p.~102]
``Rather, whether or not a competitive advantage is sustained depends upon the possibility
of competitive duplication. Following Lippman and Rumelt (1982) and Rumelt (1984), a
competitive advantage is sustained only if it continues to exist after efforts to
duplicate that advantage have ceased.''}. This review therefore keeps the two apart
throughout, because a firm can raise its margins without holding anything a rival is
unable to copy.

Where both constructs appear, the separation is visible in the survey items.
Chatterjee et al. ask about rivals directly: one of their competitive-advantage items
has respondents affirm that successful implementation of AI-based customer relationship
management ``will help an organization to win over its competitor which has not yet
implemented AI-CRM for B2B relationship management'' \parencite[p.~215]{chatterjee2021effect}\footnote{[Chatterjee et al. (2021),
p.~215] ``I think that successful implementation of AI-CRM will help an organization
to win over its competitor which has not yet implemented AI-CRM for B2B relationship
management.''}. Their performance items stay inside the firm and ask whether the same
system improves ``operational efficiency'' and speeds up
decisions \parencite[p.~215]{chatterjee2021effect}\footnote{[Chatterjee et al. (2021),
p.~215] ``I believe that successful implementation of AI-CRM will help the organization
to improve its operational efficiency.'' / ``I think AI-CRM helps in quick decision
making which helps the organization to improve its operational performance.''}.

Mehta et al. give competitive advantage a position between adoption and results: they
test it as the channel through which AI adoption reaches performance, and predict a
positive effect on performance in its own
right \parencite[p.~550]{mehta2025driven}\footnote{[Mehta et al. (2025), p.~550]
``H8. Competitive advantage has a significant and positive impact upon firm
performance.'' / ``H9. Competitive advantage has a mediating influence on adoption of
AI technologies and firm performance.''}.

Not every study that names competitive advantage builds its own instrument. Hossain
et al. measure sustained competitive advantage with items carried over from earlier
work \parencite[p.~247]{hossain2022marketing}\footnote{[Hossain et al. (2022), p.~247]
``sustained competitive advantage items were taken from Cao et al. (2019).''}. The
construct label alone says little about what a study measured; the items decide.

\section{Artificial intelligence as a general-purpose technology}
\label{sec:gpt}

Economists file AI under general-purpose technologies (GPTs), and that classification
already predicts an uneven payoff. A GPT earns the label because it makes ``new and
complementary production methods'' available, which can raise output over
time \parencite[p.~188]{czarnitzki2023artificial}\footnote{[Czarnitzki et al. (2023),
p.~188] ``The main feature of general-purpose technology is to enable new and
complementary production methods that may increase productivity over time (Bresnahan
and Trajtenberg 1995; Bresnahan et al., 2002; Brynjolfsson and Hitt 2003; Cardona
et al., 2013).''}. Those complements are the expensive part. Before output moves, a
firm has to redesign how work is organised, retrain staff, rework its software and
accumulate managerial experience, and because none of that accumulation is booked as
an asset, measured productivity dips before it
climbs \parencite[pp.~333--334]{brynjolfsson2021productivity}\footnote{[Brynjolfsson
et al. (2021), pp.~333--334] ``realizing their potential also requires large intangible
investments and a fundamental rethinking of the organization of production itself.
Firms must create new business processes, develop managerial experience, train workers,
patch software, and build other intangibles. This raises productivity measurement
issues because intangible investments are not readily tallied on a balance sheet or in
the national accounts.'' / [p.~334] ``The error in measured total factor productivity
growth therefore follows a J-curve shape, initially dipping while the investment rate in
unmeasured capital is larger than the investment rate in other types of capital, then
rising as growing intangible stocks begin to contribute to measured
production.''}. The lag is years, not quarters.

Two studies in this corpus show how differently that plays out at firm level. On German
innovation-survey panel data, Czarnitzki et al. measure a clear productivity gain from
AI use, for the yes/no measure and for the intensity measure
alike \parencite[pp.~189, 201]{czarnitzki2023artificial}\footnote{[Czarnitzki et al.
(2023), p.~189] ``We analyze cross-section as well as panel data from the German part of
the European Commission's Community Innovation Survey (CIS).'' / [p.~201] ``We found
that employing AI technologies has a positive and significant impact
on firm productivity. In particular, we showed that both the use of AI and the intensity
with which firms exploit the potential of AI significantly increase both sales and
value-added.''}. Fontanelli et al. ask a different question of French firms and find
that AI use raises the \emph{volatility} of productivity growth; the effect sits with
firms that buy AI from outside providers, and it weakens where more of the payroll is
made up of ICT engineers and
technicians \parencite[pp.~1--2]{fontanelli2025predictive}\footnote{[Fontanelli et al.
(2025), p.~1] ``This paper shifts the focus and investigates the impact of predictive AI
on the volatility of firms' productivity growth rates.'' / ``Using firm-level data from
the 2019 French ICT survey, we provide robust evidence that AI use is associated with
increased volatility.'' / [pp.~1--2] ``Our results show that heightened volatility is concentrated among AI
buyers, whereas firms that develop AI internally experience no such association.
Finally, we find that the AI-volatility link among `AI buyers' is mitigated in firms
with a higher share of ICT engineers and technicians, suggesting that AI's successful
integration requires complementary human capital.''}. Across the two results, the
reading this review takes forward is that the same AI investment can stand behind
either outcome; what differs between the firms is the complements.

\section{The resource-based view and the imitability problem}
\label{sec:rbv}

Across the studies reviewed here, the resource-based view (RBV) is the recurring reason
given for expecting AI to pay off. Barney sets four conditions on a resource before it
can hold an advantage in place: it has to be valuable, rare among competitors and
imperfectly imitable, and no substitute may do the same strategic
job \parencite[pp.~105--106]{barney1991firm}\footnote{[Barney (1991), pp.~105--106] ``To
have this potential, a firm resource must have four attributes: (a) it must be valuable,
in the sense that it exploit opportunities and/or neutralizes threats in a firm's
environment, (b) it must be rare among a firm's current and potential competition, (c) it
must be imperfectly imitable, and (d) there cannot be strategically equivalent substitutes
for this resource that are valuable but neither rare or imperfectly imitable.''
Typographical errors as in the original.}. Alwakid and Dahri restate the four conditions
as the VRIN
acronym \parencite[p.~14]{alwakid2025harnessing}\footnote{[Alwakid \& Dahri (2025), p.~14]
``Traditionally, RBV posits that firms achieve sustained competitive advantage by
strategically deploying valuable, rare, inimitable, and non-substitutable (VRIN)
resources (Barney, 1991; Wang \& Zhang, 2025).''}. That is a shorthand Barney himself comes close
to when he lists ``rareness, inimitability, non-substitutability'' among the
characteristics that can qualify a firm attribute as a source of advantage, adding that
the attribute becomes a resource only once it exploits an opportunity or neutralizes a
threat \parencite[p.~106]{barney1991firm}\footnote{[Barney (1991), p.~106] ``Firm
attributes may have the other characteristics that could qualify them as sources of
competitive advantage (e.g., rareness, inimitability, non-substitutability), but these
attributes only become resources when they exploit opportunities or neutralize threats in
a firm's environment.''}. Hossain et al. put the same argument
in terms of what a firm does with what it owns, and tie the long-run case to holdings
rivals cannot match \parencite[pp.~241--242]{hossain2022marketing}\footnote{[Hossain
et al. (2022), pp.~241--242] ``The resource-based view (RBV) portrays how a firm can
optimally use its resources, which are non-substitutable, unique and valuable for
achieving sustained competitive advantage (Wernerfelt, 1984).''}.

Purchased AI fails that test, and the RBV literature says so plainly. Krakowski et al.
note that AI costs almost nothing to reproduce and that little stands in the way of
copying it, so where AI substitutes a human capability, the RBV predicts that the
advantage resting on that capability will
erode \parencite[p.~1426]{krakowski2023artificial}\footnote{[Krakowski et al.
(2023), p.~1426] ``When AI substitutes humans' cognitive capabilities, the RBV expects
the advantage that these capabilities provided to erode (Peteraf \& Bergen, 2003). This
is because, as a technology resource, AI has close to zero marginal reproduction costs
and few imitation barriers (Brynjolfsson \& McAfee, 2014, p.~31).''}. Their own evidence
relocates the advantage rather than denying it: what resists imitation is the combination
a firm builds, not either side of
it \parencite[p.~1446]{krakowski2023artificial}\footnote{[Krakowski et al. (2023),
p.~1446] ``Our findings show that humans remain sources of sustainable advantage by
combining and integrating their capabilities with machine capabilities into uniquely
complementary bundles (Argyres \& Zenger, 2012).'' / ``the locus of sustainable advantage
is neither human nor machine capabilities, but what humans do with these capabilities.''}.

Two consequences follow for this review. The RBV speaks to competitive advantage and not
to performance, so a study reporting higher margins has not thereby shown that anything
is defensible against rivals; the separation of the two constructs in
Section~\ref{sec:constructs} is a theoretical requirement, not only a coding convention.
And if the resource itself is imitable, the variation worth studying lies in what
surrounds it, which is why the conditions under which AI pays off are treated here as the
object of analysis rather than as controls.
% Barney (1991) wird ueber literature/barney1991firm.txt zitiert (verifiziertes
% Text-Sidecar zum Bildscan, Seitenmarker = Journalseiten). Alle Passagen am
% Seitenbild gegengelesen: S. 102 (Definitionen, in Abschnitt 2.1) sowie
% S. 105-106 und S. 106 (vier Attribute, hier).

\section{Automation and augmentation}
\label{sec:autoaug}

The same investment can be deployed in two ways, and the strategy literature keeps them
apart. Raisch and Krakowski define automation as machines taking a task over from a
human, and augmentation as humans and machines working on a task
together \parencite[p.~192]{raisch2021artificial}\footnote{[Raisch \& Krakowski (2021),
p.~192] ``Whereas automation implies that machines take over a human task, augmentation
means that humans collaborate closely with machines to perform a task.''}. Their argument
against the popular advice to prefer augmentation is that the two cannot be held apart in
practice. For any single task at any single moment the firm picks one approach, and that
choice makes the two contradictory; widen the frame to several tasks over several years
and each turns out to produce the other, which makes them
interdependent \parencite[pp.~195--197]{raisch2021artificial}\footnote{[Raisch \&
Krakowski (2021), p.~195] ``Automation and augmentation are contradictory, because
organizations choose either one or the other approach to address a given task at a
specific point in time.'' / ``If we increase our analysis's temporal scale (from one
point in time to its evolution over time) and spatial scale (from one to multiple
tasks), we comprehend that, in the management domain, the two AI applications are not
only contradictory but also interdependent.'' / [p.~196] ``Consequently, augmentation
may enable a transition to automation over time.'' / [p.~197] ``Automation in one task
`spills over,' enabling adjacent tasks' augmentation.''}.

Corpus studies pick the distinction up as a design choice with different returns. Zebec
and Indihar {\v S}temberger treat the two as the use cases through which AI adoption
reaches performance, automation paying off in efficiency and augmentation in what people
and machines manage
jointly \parencite[p.~7]{zebec2024creating}\footnote{[Zebec \& Indihar {\v S}temberger
(2024), p.~7] ``Automation involves machines taking over tasks from humans, leading to
performance gains via resource, cost and information-processing efficiencies. In
contrast, augmentation involves humans collaborating with machines, enhancing both human
and machine skills and productivity.''}. Their model then returns no significant direct
effect of process automation on process performance once tasks become complex, and the
recommendation they draw from it is a boundary condition: automate what is routine,
augment what is
ambiguous \parencite[pp.~16--17]{zebec2024creating}\footnote{[Zebec \& Indihar
{\v S}temberger (2024), pp.~16--17] ``The results show that BPA has no significant direct
effect on BPP for more complex tasks in terms of efficiency of execution or scalability.''
/ ``In summary, automate routine, well-structured tasks and augment complex, ambiguous
ones.''}. Where the line falls is itself unstable. In the fragrance industry the machine
took over part of the perfumers' problem solving and could not touch the part that
depends on smelling a scent and judging how people will react to
it \parencite[p.~1445]{krakowski2023artificial}\footnote{[Krakowski et al. (2023),
p.~1445] ``While machines substitute perfumers regarding certain problem solving
activities, they are unable to match the perfumers' unique ability to smell fragrances
and predict the human emotions they trigger. This example illustrates that AI is likely
to substitute some, but not all of complex business tasks' activities.''}.

\section{Situated AI}
\label{sec:situated}

Kemp starts from the puzzle the two previous sections produce: AI is available to
everyone, so buying it should not yield an advantage, yet firms plainly differ in what
they get out of it. His answer is situated AI, which he defines as AI hemmed in by the
firm's own experience, structure and
relationships \parencite[p.~618]{kemp2024competitive}\footnote{[Kemp (2024), p.~618]
``The present paper begins to resolve this puzzle by developing a theory of situated
AI---AI whose agency is circumscribed in a firm's experiential, structural, and
relational systems.''}. Three properties of the technology stand in the way, and the
first of them is the imitability problem from Section~\ref{sec:rbv} under another name:
AI is generic, explicit and myopic. Being generic puts AI on a level with general human
capital, which competitors can acquire and put to work much as the adopter
does \parencite[p.~620]{kemp2024competitive}\footnote{[Kemp (2024), p.~620] ``The generic,
explicit, and myopic nature of AI all endanger its firm-level benefits, but are uniquely
difficult to surmount because AI is both a machine and agentic (Smith \& Lewis, 2011).
For example, AI's generic nature is akin to general human capital. General human capital
cannot typically underly interfirm advantages because a firm's competitors can usually
acquire and deploy that knowledge in ways that closely mimic corresponding actions of the
focal firm.''}.

What a firm does control is the situating. Kemp names three activities: grounding, which
picks the experiences a firm's own AI learns from; bounding, which works on the
experiences feeding a competitor's AI; and recasting, which keeps reworking the
algorithms and the routines around them so the AI stays aligned with the rest of the
firm's activities \parencite[pp.~618--619]{kemp2024competitive}\footnote{[Kemp (2024),
pp.~618--619] ``Grounding involves orchestrating which experiences one's AI will be
allowed to learn from across the organization. Bounding involves efforts to shape the
experiences anchoring a competitor's AI. Recasting involves orchestrating the continual
adaptation of algorithms and their surrounding routines to enhance AI's alignment with
interdependent activities in a firm.''}. The vocabulary is conditional throughout. None
of the three is a property of the technology, and each is something a firm may or may not
do.

Empirical work in the corpus starts from the same reading. Shi et al. take the mixed
results of earlier work as a sign that what AI yields depends on intangible complements,
on governance, and on managerial capability bound to its context, and they test that
claim on Chinese listed firms, with data assets as the channel and managerial capability
as the moderator \parencite[p.~1]{shi2025value}\footnote{[Shi et al. (2025), p.~1] ``Drawing on a
comprehensive dataset of Chinese listed companies spanning 2010--2022, this study
investigates how AI adoption influences enterprise value through the mediating role of
data assets and the moderating effect of managerial capability.'' / ``These mixed
outcomes suggest that the value-creating capacity of AI depends on more than technological
sophistication alone; it hinges on complementary intangible resources, effective
governance structures, and context-specific managerial capabilities---factors that current
empirical research has only partially explored.''}. Kemp's paper is conceptual, so what it
offers are propositions rather than findings. They are the propositions this review is
positioned to examine, because the conditions Kemp derives are the conditions the coded
studies report.
